\section{The return map and uniform estimates}
\label{sec:passage}

The two contact boxes and their regular connections determine a return map on
one transverse section. We first establish its itinerary and the reversible
passage law, then derive the cancellation and the value estimates used in the
fixed-point count.

\emph{Conventions.} A family is fixed throughout. All domains and bounds are
chosen on compact reductions that are common to the whole family and do not
depend on how the parameters approach the reference; derivative orders are
prescribed in advance. Derivatives in the state are taken at fixed parameters,
and derivatives and multiplicities refer to the transverse coordinate $q$. A
function of $\sqrt q$ that extends smoothly through $\sqrt q=0$ is called a
\emph{root profile}; smoothness of a root profile does not imply smoothness in
$q$. The letters $K$ and $C_j$ denote constants that may change from line to
line, and letters for smooth remainders, such as $L,M,N,S,E$, are local to the
statement or proof in which they appear.

\subsection{Geometry and a common first return}
\label{subsec:observed-controls}

\begin{definition}[Admissible family]\label{def:admissible-family}
A finite-dimensional smooth family of continuously matched planar fields,
parametrized by $\lambda$, is \emph{admissible} if the following hold near
its reference parameter $\lambda_0$.
\begin{enumerate}
\item The branch fields, switching curves, marked contacts and transverse
sections are jointly smooth. In each of two disjoint contact boxes,
\begin{equation}\label{eq:matched}
 X^+_\lambda=X^-_\lambda\quad\hbox{on }\{h_\lambda=0\},
 \qquad X^+_\lambda-X^-_\lambda=h_\lambda Z_\lambda.
\end{equation}
The switching function $h_\lambda$ has nonzero differential, the inactive
field $X^-_\lambda$ does not vanish, and its marked orbit has a persistent
quadratic contact: at each marked contact,
\[
 X^-_\lambda h_\lambda=0,\qquad
 (X^-_\lambda)^2h_\lambda\le-\kappa_0<0
\]
for one common $\kappa_0>0$. The active region is $\{h_\lambda>0\}$, where
$X^+_\lambda$ is used; near each marked contact the marked orbit lies in
$\{h_\lambda\le0\}$.

\item Each contact box has an incoming and an outgoing transverse section
(its ports), with closures in the interior of the box, and consecutive boxes
are joined by orbit arcs of $X^-_\lambda$ along which $h_\lambda<0$. On a
common neighbourhood of states and parameters, an orbit that enters a box with
positive penetration (the local input $z$ defined below) crosses the seam into
the active region, stays there, and leaves through the seam before reaching
the outgoing port; an orbit with nonpositive penetration touches the seam at
most once and otherwise stays in $\{h_\lambda<0\}$. On one fixed incoming
interval $I=(-q_0,q_0)$, including its grazing thresholds, every orbit follows
this itinerary, meets no seam along the connecting arcs, and does not return
to the incoming section before completing one turn. All intermediate port
images lie in common open domains.
The local penetration coordinates are positive towards the active regions.
The regular transition from the first outgoing port to the second incoming
port is increasing in these coordinates, with a common positive derivative
margin. Thus both active sides agree after transport to the incoming section.
\item At $\lambda_0$ there is a compact embedded nonstationary periodic orbit
$\mathcal C$ through the two marked contacts. Its return-map controls, passage
strengths and closing curvature, defined in~\eqref{eq:net-strength}
and~\eqref{eq:regular-jets} below, satisfy
\[
 a=b=\delta=0,\qquad A_1=-A_2\ne0,\qquad c_0\ne0.
\]
\end{enumerate}
\end{definition}

We transport the inactive first integrals along the regular connections to
one smooth transverse coordinate $q$, with contact levels $\tau_1=0$ and
$\tau_2=\delta$; Proposition~\ref{prop:field-recognition} below shows that
this is possible for nonidentity connecting maps. In the $\nu$th box put
$z=q-\tau_\nu$. An inactive flow box gives the seam $z=s^2\gamma(s,\lambda)$
with $\gamma>0$; the signed clock coordinate $t=s\sqrt{\gamma(s,\lambda)}$
leaves $z$ unchanged and normalizes the seam. After a positive time change,
matching and smooth division give
\begin{equation}\label{eq:passage:local-geometry}
 X^-=\partial_t,\qquad \Sigma=\{z=t^2\},\qquad
 X^+=(1+h\alpha)\partial_t+h\mu\partial_z,
 \qquad h=z-t^2.
\end{equation}
The functions $\alpha,\mu$ are jointly smooth and the active clock is
positive on a common smaller box. This normalization preserves oriented
orbits and section maps, so it does not change the return map.
Figure~\ref{fig:actual-return-geometry} shows the reference itinerary and the
grazing thresholds on the incoming section.
% SCHEMATIC: no numerical trajectory, metric scale, or stability is encoded.
% Top: the reference cycle at delta=0, oriented by time.
% Bottom: a perturbed incoming-state axis for delta>0, not phase-plane distance.
\begin{figure}[!htbp]
  \centering
  \begin{tikzpicture}[x=1cm,y=1cm,font=\small,
      orbit/.style={line width=.85pt},
      flow/.style={line width=.85pt,-{Stealth[length=1.8mm]}},
      seam/.style={line width=.65pt,densely dashed,black!65},
      box/.style={line width=.5pt,densely dotted,black!60},
      port/.style={line width=.75pt},
      axis/.style={line width=.55pt,-{Stealth[length=1.8mm]}}]
    \begin{scope}[yshift=2.9cm]
      \node[anchor=west] at (0,3.85)
        {(a) Reference itinerary, $\delta=0$};

      \draw[box] (1.5,1.45) rectangle (4.1,3.12);
      \draw[box] (6.0,1.45) rectangle (8.6,3.12);
      \node at (2.8,3.35) {$U_1:\ \Gamma_1$};
      \node at (7.3,3.35) {$U_2:\ \Gamma_2$};
      \draw[seam] (1.75,2.71625) parabola bend (2.8,2) (3.85,2.71625);
      \draw[seam] (6.25,2.71625) parabola bend (7.3,2) (8.35,2.71625);
      \node[anchor=west] at (3.75,2.95) {$\Sigma_1$};
      \node[anchor=west] at (8.25,2.95) {$\Sigma_2$};

      \draw[orbit] (.6,2) -- (4.1,2);
      \draw[flow] (4.1,2)
        .. controls (4.4,2) and (4.55,2.25) .. (5.05,2.25)
        .. controls (5.55,2.25) and (5.7,2) .. (6.0,2);
      \draw[orbit] (6.0,2) -- (8.6,2);
      \draw[flow] (8.6,2) .. controls (10.1,2) and (10.1,.3) .. (8.6,.3);
      \draw[orbit] (8.6,.3) -- (1.0,.3)
        .. controls (-.35,.3) and (-.35,2) .. (.6,2);
      \draw[flow] (.85,2) -- (1.25,2);
      \draw[flow] (3.2,2) -- (3.62,2);
      \draw[flow] (7.65,2) -- (8.07,2);
      \draw[flow] (6.0,.3) -- (5.15,.3);
      \node[align=center,text width=1.65cm] at (5.05,2.83)
        {regular\\connection};
      \node at (5.0,.7) {closing leg $R$};
      \node at (5.05,1.35) {$\mathcal C$};

      \draw[port] (.6,1.55) -- (.6,2.65);
      \draw[axis] (.2,2.12) -- (.2,2.62) node[above] {$q$};
      \node[anchor=east] at (.42,1.5) {$\mathcal S$};
      \draw[port] (1.75,1.87) -- (1.75,2.13);
      \draw[port] (3.85,1.87) -- (3.85,2.13);
      \draw[port] (6.25,1.87) -- (6.25,2.13);
      \draw[port] (8.35,1.87) -- (8.35,2.13);
      \node[below] at (1.75,1.4) {$\Pi_1^{\rm in}$};
      \node[below] at (3.85,1.4) {$\Pi_1^{\rm out}$};
      \node[below] at (6.25,1.4) {$\Pi_2^{\rm in}$};
      \node[below] at (8.35,1.4) {$\Pi_2^{\rm out}$};
      \draw[fill=white,line width=.8pt] (2.8,2) circle[radius=1.8pt];
      \draw[fill=white,line width=.8pt] (7.3,2) circle[radius=1.8pt];
      \node[below] at (2.8,1.93) {$p_1$};
      \node[below] at (7.3,1.93) {$p_2$};
    \end{scope}

    \node[anchor=west] at (0,2.2)
      {(b) Incoming-state activation, $\delta>0$};
    \draw[axis] (.5,1.05) -- (10.25,1.05) node[right] {$q$};
    \draw[port] (3.0,.94) -- (3.0,1.16);
    \draw[port] (6.5,.94) -- (6.5,1.16);
    \node[below] at (3.0,.9) {$0$};
    \node[below] at (6.5,.9) {$a_\delta=\Gamma_1^{-1}(\delta)$};
    \node at (1.5,1.5) {no excursion};
    \node at (4.75,1.5) {first only};
    \node at (8.55,1.5) {both excursions};
    \node at (5.25,.12) {second local input: $\Gamma_1(q)-\delta$};
  \end{tikzpicture}
  \caption{Schematic return geometry. Top: the reference itinerary at
    $\delta=0$, with dotted contact boxes $U_1,U_2$, dashed switching seams
    $\Sigma_1,\Sigma_2$, transverse ports and the return section
    $\mathcal S$. Arrows on the cycle show the direction of time; the
    section arrow orients $q$. Bottom:
    activation on the incoming state axis for $\delta>0$. The second
    passage is active when $\Gamma_1(q)-\delta>0$, so its threshold is
    $a_\delta$. For $\delta<0$, $a_\delta=\delta<0$ lies to the left of
    $0$, and only the second excursion occurs between them.}
  \label{fig:actual-return-geometry}
\end{figure}


\begin{lemma}[A common one-turn return tube]\label{lem:common-tube}
Let a compact embedded nonstationary reference cycle admit fixed tubular
coordinates $(\theta,w)$, with $\theta\in\R/\mathbb Z$, reference $w=0$
and section $\theta=0$. Suppose a finite smooth family of matched fields
has common local Lipschitz bounds on the tube and differs from the
reference field at $\lambda_0$ by $O(|\lambda-\lambda_0|)$. Fix contact cores and transverse ports with
strict margins, and strictly inactive connecting arcs. In every core,
assume that orbits with positive penetration enter the active region, stay
there and leave through the seam inside the core, and that the other orbits
touch the seam at most once.
If the angular clock remains positive and these local and port margins
persist, one can choose a common input interval and a parameter neighbourhood
on which the chronological composition is the first return map. It makes one
whole turn, with no further seam crossing or earlier section hit.
\end{lemma}

\begin{proof}
The matched form $X^-_\lambda+(h_\lambda)_+Z_\lambda$ has common
Lipschitz bounds. Compactness gives $\dot\theta\ge v_*>0$ on a smaller
tube; division by this clock and Gronwall comparison with $w=0$ give,
up to a first exit,
\[
 \begin{aligned}
 \left|\frac{dw}{d\theta}\right|&\le L|w|+K|\lambda-\lambda_0|,\\
 |w(\theta)|&\le e^L(|w(0)|+K|\lambda-\lambda_0|),\quad0\le\theta\le1.
 \end{aligned}
\]
Choose one input interval and parameter neighbourhood keeping this bound
strictly inside the tube, core and port margins. First-exit continuation
gives a full turn; the assumed core passages and inactive arcs exclude
further seam crossings, including for tangent inputs. The ports join on common
domains, and strict angular advance makes $\theta=1$ the next section hit.
The return image need not lie in the input interval.
\end{proof}

Lemma~\ref{lem:common-tube} gives a geometric criterion for the first-return
itinerary in item~2 of Definition~\ref{def:admissible-family}; Section~\ref{sec:sharpness} applies it
to the constructed families.

\begin{remark}[Cycles and fixed points]\label{rem:cycle-fixed-point}
After a common reduction, every periodic orbit sufficiently close to
$\mathcal C$ in the return tube meets the transverse section in $I$.
The first return is increasing and has no nontrivial finite-period orbit
whose iterates remain in the common local domain: if $P(q)>q$ or $P(q)<q$,
its successive iterates are strictly ordered. Thus these periodic orbits
correspond to fixed points in $I$, and isolation of the fixed point is
equivalent to isolation of the cycle in this tube.
\end{remark}

\subsection{Reversible passages and cancellation}
\label{sec:neutral-reference}\label{sec:center-multiplicity}

An active leaf pairs its two seam intersections by a decreasing
involution. We use it to identify the entry and exit of the passage.
The resulting relation between its first two coefficients determines
cancellation in the two-passage composition.

\begin{proposition}[Passage map near a quadratic grazing]
\label{prop:quadratic-passage}\label{prop:passage:quadratic}
For a local passage with the quadratic geometry~\eqref{eq:passage:local-geometry}
and fixed local transverse ports, there is a jointly smooth
signed companion $g(t,\lambda)$ on a two-sided interval such that
\begin{equation}\label{eq:passage:reversal}
 g(0,\lambda)=0,\qquad g_t(0,\lambda)=1,\qquad
 \iota g\iota=g^{-1},\qquad \iota(t)=-t.
\end{equation}
The passage in the inactive coordinate is
\begin{equation}\label{eq:passage:physical}
 \Gamma(q,\lambda)=q\quad(q\leq0),\qquad
 \Gamma(r^2,\lambda)=g(r,\lambda)^2\quad(r>0).
\end{equation}
Its first entry is $t=-r$ and its first moving exit is $g(r,\lambda)$.
On a common domain,
\begin{equation}\label{eq:passage:quadratic-jet}
 \Gamma(q,\lambda)
 =q+C_1(\lambda)q_+^{3/2}
       +\frac34C_1(\lambda)^2q_+^2
       +q_+^{5/2}L(\sqrt{q_+},\lambda),
\end{equation}
where $L$ is the restriction of a jointly smooth function, and
\begin{equation}\label{eq:passage:strength-jet}
 C_1(\lambda)=\frac43\mu(0,0,\lambda).
\end{equation}
The passage is an orientation-preserving $C^1$ local diffeomorphism, and the
companion, its inverse and the root profiles satisfy uniform bounds of every
prescribed order.
\end{proposition}

\begin{proof}
Divide the active field by its positive $t$ component. Let
$K(t,z,\lambda)$ be the $z$ coordinate where its orbit meets $t=0$.
Smooth finite-time ODE dependence gives $K(0,z,\lambda)=z$ and $K_z>0$
on a common smaller box. Since the active orbit equation has zero
$z$ component on the seam, $K_t(t,t^2,\lambda)=0$. Hence the seam
restriction $\widehat K(t,\lambda)=K(t,t^2,\lambda)$ satisfies
\[
 \begin{gathered}
 \widehat K_t=2tE(t,\lambda),\qquad
 E(t,\lambda)=K_z(t,t^2,\lambda)>0,\\
 \widehat K=t^2M(t,\lambda),\qquad
 M(t,\lambda)=2\int_0^1 sE(st,\lambda)\,ds.
 \end{gathered}
\]
Here $E(0,\lambda)=M(0,\lambda)=1$, so
$\mathcal N(t,\lambda)=t\sqrt{M(t,\lambda)}$ is a smooth increasing
local diffeomorphism. The active leaf pairs its seam intersections by
\[
 J=\mathcal N^{-1}\iota\mathcal N,\qquad
 J(0)=0,\quad J'(0)=-1,\quad J^2=\mathrm{id}.
\]

For input $q=r^2>0$, the inactive orbit first reaches the seam at $-r$.
Its active leaf has value $\widehat K(-r)>0$ and second intersection
$e=J(-r)>0$. The seam restriction decreases for $t<0$ and increases
for $t>0$; since $K_z>0$, the leaf $K=\widehat K(-r)$ lies strictly above the seam
precisely between $-r$ and $e$. The positive clock makes them the first
entry and first exit. In the normalized clock, $dh/dt=-2t$ on the seam, so the entry and exit
derivatives are exactly $2r$ and $-2e$. In the original physical time
they are multiplied by the positive clock factor. After exit the inactive orbit has $z=e^2$
and lies below the seam; port localization excludes further seam crossings.
Inputs $q<0$ stay below it and $q=0$ only touches it, giving the identity
passage on $q\le0$.

Set $g=J\iota$. Then $J^2=\iota^2=\mathrm{id}$ gives
$\iota g\iota=g^{-1}$ and $g'(0)=1$, while $e=g(r)$ proves
\eqref{eq:passage:physical}. For negative $t$, reversal gives
\begin{equation}\label{eq:passage:negative-companion}
 g(t,\lambda)=-\sqrt{\Gamma^{-1}(t^2,\lambda)}\qquad(t<0).
\end{equation}
This signed continuation records backward active travel; the forward
passage on negative $q$ is the identity.

Writing $g(t)=t+\rho t^2+\sigma t^3+O(t^4)$ and comparing cubic
coefficients in $(g\iota)^2=\mathrm{id}$ gives $\sigma=\rho^2$, the
classical jet relation for reversible maps~\cite{ofarrell2008,novaessilva2021}.
Taylor's integral remainder therefore gives
\[
 g(t)=t+\rho t^2+\rho^2t^3+t^4G(t,\lambda)
\]
with $G$ jointly smooth. Squaring at $t=\sqrt q>0$ gives
$C_1=2\rho$ and $C_2=3\rho^2=3C_1^2/4$, proving
\eqref{eq:passage:quadratic-jet}. Its parameter derivatives also give
joint $C^1$ regularity and $\Gamma_q\to1$ at zero.

For~\eqref{eq:passage:strength-jet}, note that $dz/dt=f(t,z)=h\mu/(1+h\alpha)$
and $K_t=-fK_z$. On the seam $f=0$ and $f_z=\mu$, so $E'(0)=-\mu(0,0)$.
Thus $\mathcal N(t)=t-\mu(0,0)t^2/3+O(t^3)$, and
$\mathcal N^{-1}\iota\mathcal N\iota$ has quadratic coefficient $2\mu(0,0)/3$;
squaring gives~\eqref{eq:passage:strength-jet}. Compactness, smooth dependence
on initial data and a common shrinking give uniform positivity and the
prescribed derivative bounds.
\end{proof}

Since $\partial_q=(2r)^{-1}\partial_r$, the $j$th state derivative of
$q^{5/2}L(\sqrt q,\lambda)$ is bounded by $C_jq^{5/2-j}$ on $q>0$, for every
fixed $j$ and after finitely many parameter derivatives. Its zero extension to
$q\le0$ is jointly $C^2$ in state and parameters.

Write $\Gamma_1,\Gamma_2$ for the two normalized passages.
The inactive connections preserve $q$, so chronological continuation gives
\begin{equation}\label{eq:passage:two-contact-map}
 \begin{aligned}
 Q_{m,\delta}(q,p)&=\delta+
 \Gamma_2\bigl(\Gamma_1(q,p,m,\delta)-\delta,p,m,\delta\bigr),\\
 P_{m,\delta}&=R_{p,m,\delta}\circ Q_{m,\delta},
 \end{aligned}
\end{equation}
where $R$ is the smooth orientation-preserving transition along the rest of
the cycle, the \emph{closing map}. Write $D=P-q$ and
$A_\nu(\lambda)=C_{1,\nu}(\lambda)$, the \emph{passage strengths}. When the
controls are independent we write $p=(a,b,\xi)$ and
$\lambda=(a,b,m,\delta,\xi)$, with $\xi$ in a fixed compact set. Throughout,
$m$ denotes the strength sum:
\begin{equation}\label{eq:net-strength}
 A_1(p,m,\delta)+A_2(p,m,\delta)=m,\qquad A(p)=A_1(p,0,0).
\end{equation}
The individual strengths and all higher profiles may depend on every
parameter. After reduction, $A$ has fixed sign with
$|A(p)|\ge A_*>0$, and the two passage strengths have opposite signs. The jets of the closing map are
\begin{equation}\label{eq:regular-jets}
 a=R(0),\qquad b=R'(0)-1,\qquad c=\tfrac12R''(0).
\end{equation}
In particular, the aligned multiplier is $k(p)=R'_{p,0,0}(0)=1+b>0$.
Taylor division gives
\begin{equation}\label{eq:quartic-closing-division}
 R(y)=a+(1+b)y+c(\lambda)y^2+y^3S(y,\lambda).
\end{equation}
At the reference $a=b=m=\delta=0$. The aligned curvature
$c_0(p)=\tfrac12R''_{p,0,0}(0)$ has fixed nonzero sign and a uniform
strict margin; after shrinking, $R''$ retains that sign at every relevant
input. Its value may vary with the parameters. The attainment results use
the condition
\begin{equation}\label{eq:rank}
 \operatorname{rank}D_\lambda(a,b,m,\delta)=4
 \quad\hbox{at the reference field},
\end{equation}
which makes $(a,b,m,\delta,\xi)$ local parameter coordinates.

\begin{lemma}[Cancellation at aligned thresholds]
\label{lem:passage:finite-inverse}
On $\delta=0$, the composition has the root-profile form
\begin{equation}\label{eq:passage:aligned-jet}
 Q_{m,0}(q)=q+mq_+^{3/2}+\frac34m^2q_+^2
                  +q_+^{5/2}L_0(\sqrt{q_+},a,b,m,\xi).
\end{equation}
Consequently, on the whole slice $m=\delta=0$,
\begin{equation}\label{eq:passage:finite-cancellation}
 Q_{0,0}(q)=q+q_+^{5/2}L_0(\sqrt{q_+},a,b,0,\xi)
\end{equation}
is jointly $C^2$. Only the leading strengths have to cancel.
\end{lemma}

\begin{proof}
On a common smaller domain positive inputs and intermediate images stay
positive. The cubic companion jets and $A_1+A_2=m$ give
\[
 g_2g_1(t)=t+\frac m2t^2+\frac{m^2}{4}t^3
                      +t^4G_0(t,a,b,m,\xi),
\]
with $G_0$ jointly smooth. Squaring proves
\eqref{eq:passage:aligned-jet}; for $q\le0$ both passages are identity.
At $m=0$ the remaining $q_+^{5/2}$ profile has the joint $C^2$
extension proved above.
\end{proof}

The grazing thresholds on the incoming section are
\begin{equation}\label{eq:passage:true-thresholds}
 0\quad\hbox{and}\quad a_\delta=\Gamma_1^{-1}(\delta,\lambda).
\end{equation}
Here $a_\delta$ labels the second threshold; it does not denote a
parameter derivative. The second local input is
$I_2=\Gamma_1(q,\lambda)-\delta$.
The increasing first passage has derivative near one, so its inverse
exists on a common interval and
\begin{equation}\label{eq:passage:threshold-formula}
 a_\delta=\delta\quad(\delta\le0),\qquad
 a_\delta=\bigl[g_1^{-1}(\sqrt\delta,\lambda)\bigr]^2
                                                \quad(\delta>0).
\end{equation}
Thus $a_\delta$ has the sign of $\delta$, $|a_\delta|\le2|\delta|$,
and for positive separation
$a_\delta=\delta-C_{1,1}(\lambda)\delta^{3/2}+O(\delta^2)$ with a
smooth positive-root profile. It is jointly $C^1$; as a function of the
independent separation variable $\delta$, it is not $C^2$ at $\delta=0$
when $C_{1,1}\ne0$.
For $\delta>0$, the intervals $q<0$, $0<q<a_\delta$, $q>a_\delta$
have respectively no excursion, only the first, and both excursions.
For $\delta<0$, the corresponding intervals are $q<\delta$,
$\delta<q<0$, $q>0$, with only the second excursion in the middle.
At either endpoint its excursion becomes a tangency touch. At alignment,
negative inputs are inactive, zero touches both contacts and positive
inputs make both excursions.

On the aligned cancelling slice, the closing expansion gives
\begin{equation}\label{eq:center-displacement}
 D_{0,0}(q,p)=a+bq+c_0(p)q^2
       +q_+^{5/2}M(\sqrt{q_+},p)+q^3S(q,p),
 \qquad c_0(p)=\tfrac12R_{p,0,0}''(0)\ne0.
\end{equation}
The root remainder is jointly $C^2$, with bounds
$|\partial_q^j\partial_p^\beta(q^{5/2}M)|\le C_{j,\beta}q^{5/2-j}$
for $j\le2$. At $a=b=0$ the reference displacement consequently has
$D(0)=D_q(0)=0$ and $D_{qq}(0)=2c_0\ne0$.
The return map is jointly $C^1$ in state and parameters. Since both strengths
are nonzero, $P_\lambda$ is $C^2$ at both thresholds if and only if
$m(\lambda)=\delta(\lambda)=0$; otherwise a nonzero coefficient of order
$q^{3/2}$ survives at one of the thresholds. For the counting bounds, fixed points at thresholds are counted as distinct
zeros; no smooth multiplicity is used at a non-$C^2$ threshold.

A smooth increasing change of transverse coordinate $x=\phi(q)$, $\phi(0)=0$,
$\phi'(0)=s>0$, changes the reference coefficients to $\widetilde c_0=c_0/s$
and $\widetilde A_\nu=A_\nu(\lambda_0)/\sqrt{s}$; the terms involving $\phi''(0)$
cancel at a fixed point with unit multiplier. In particular $c_0/A_1(\lambda_0)^2$
does not depend on the coordinate. The odd root lift $t\sqrt{\phi(t^2)/t^2}$
preserves reversal, and nonlinear terms of $\phi$ first affect the order
$5/2$, so the coefficient law of Proposition~\ref{prop:quadratic-passage} is
preserved. These invariants let us check the sign conditions in any such
section coordinate and compare the constants in the weighted limits.

\subsection{Regular connections and auxiliary parameters}

We now show how regular connecting maps lead to this form, and how families
with dependent controls are handled.

\begin{proposition}[Return map through regular connections]
\label{prop:field-recognition}\label{prop:passage:field-recognition}
In the common geometry above, let $S_1,S_2$ be the regular transitions in
positive penetration coordinates, with all intermediate images in their
common domains. They are smooth and increasing, with common positive
derivative margins. If the second local level is $\delta_2$, then
\begin{equation}\label{eq:passage:unreduced-return}
 P=S_2\circ\Pi_2\circ S_1\circ\Gamma_1,
 \qquad\Pi_2(z)=\delta_2+\Gamma_{2,\mathrm{loc}}(z-\delta_2)
\end{equation}
has the transported form
\begin{equation}\label{eq:passage:recognized-return}
 P(q)=R\{\delta+\Gamma_2(\Gamma_1(q)-\delta)\},
 \qquad\delta=S_1^{-1}(\delta_2),\quad R=S_2\circ S_1.
\end{equation}
The second passage retains reversal and the coefficient law. The four
controls are smooth functions of $\lambda$, and under~\eqref{eq:rank} they
give a chart $(a,b,m,\delta,\xi)$ with $m=C_{1,1}+C_{1,2}$.
\end{proposition}

\begin{proof}
The first transition is increasing by admissibility. The inactive reference
monodromy preserves transverse orientation, so $S_2'>0$ at the reference as
well; smooth dependence gives common positive derivative margins after
reduction.
Set $T(x)=S_1(\delta+x)-\delta_2$ and
$\Gamma_2=T^{-1}\Gamma_{2,\mathrm{loc}}T$.
Then $T(0)=0$, $T'>0$, and
$S_1^{-1}\Pi_2S_1(y)=\delta+\Gamma_2(y-\delta)$, proving the formula.
The signed root lift is
\[
 H(t)=t\sqrt{T(t^2)/t^2},\qquad
 T(t^2)/t^2=\int_0^1T'(st^2)\,ds,\qquad
 g_2=H^{-1}g_{2,\mathrm{loc}}H.
\]
The quotient is smooth, positive and even. Thus $H$ is odd, increasing,
with a smooth bounded inverse on a common root interval; conjugation
preserves reversal. Squaring for $t>0$ gives the transported passage.
For $\ell=T'(0)$ its first coefficients are
$C_{1,2}=\sqrt\ell\,C_{1,2,\mathrm{loc}}$ and
$C_{2,2}=\ell C_{2,2,\mathrm{loc}}$, so $C_{2,2}=3C_{1,2}^2/4$.
Nonlinear terms in $T$ first enter the order-$5/2$ remainder.
The transported first integral is $S_1^{-1}$ of the old one; its level,
first entry and moving first exit follow the same trajectories.
The common itinerary therefore gives the chronological first return.
Smooth parameter dependence gives the four control functions, and
under~\eqref{eq:rank} the submersion theorem gives the asserted chart.
\end{proof}

Within an admissible family, finite-time dependence and the implicit
function theorem for hitting times give smooth regular transitions with
common clock, derivative and domain margins. The opposite strength signs
and nonzero closing curvature persist in a neighbourhood; cancellation is
imposed only at the reference.

For dependent controls we introduce independent auxiliary parameters for
the estimates.

\begin{lemma}[Independent auxiliary strength and separation]
\label{lem:auxiliary-controls}
Let a jointly smooth finite-parameter family have the normalized return
\[
 \begin{gathered}
 P_\lambda(q)=R_\lambda\bigl(\widehat\delta(\lambda)+
 \Gamma_{2,\lambda}(\Gamma_{1,\lambda}(q)
                              -\widehat\delta(\lambda))\bigr),\\
 \widehat m(\lambda)=C_{1,1}(\lambda)+C_{1,2}(\lambda),
 \end{gathered}
\]
with the signed companions of Proposition~\ref{prop:quadratic-passage}.
Near a reference with $\widehat m=\widehat\delta=0$, there is a jointly
smooth root-profile family with independent auxiliary parameters
$(m,\delta)$, the same coefficient law and the identity
$C_{1,1}+\widetilde C_{1,2}=m$ for all auxiliary $(m,\delta)$,
whose return satisfies
\[
 \widetilde P_{\lambda,\widehat m(\lambda),
                       \widehat\delta(\lambda)}=P_\lambda.
\]
The closing map remains $R_\lambda$, and all profiles satisfy uniform
bounds.
\end{lemma}

\begin{proof}
Put $\iota(t)=-t$, $J_2=g_{2,\lambda}\iota$ and
$\mathcal N_2(t,\lambda)=(t-J_2(t,\lambda))/2$.
Since $J_2^2=\id$, $\mathcal N_2J_2=\iota\mathcal N_2$ and
$\mathcal N_2'(0)=1$, so $\mathcal N_2$ is a jointly smooth increasing
local normalizer. For $s=m-\widehat m(\lambda)$ set
\[
 \mathcal N_s(t,\lambda)=\mathcal N_2(t,\lambda)-\frac s4t^2,
 \qquad \widetilde g_2=\mathcal N_s^{-1}\iota\mathcal N_s\iota.
\]
A common shrinking keeps $\mathcal N_s'>0$, and the parameter inverse
function theorem gives its joint smooth inverse. Thus $\widetilde g_2$
is tangent to the identity, satisfies
$\iota\widetilde g_2\iota=\widetilde g_2^{-1}$, and equals $g_{2,\lambda}$
when $s=0$. If $\mathcal N_2=t+d_2t^2+O(t^3)$, the quadratic coefficient
of $\mathcal N_2^{-1}\iota\mathcal N_2\iota$ is $-2d_2$.
The change therefore raises the quadratic coefficient of $\widetilde g_2$
by $s/2$.
Define $\widetilde\Gamma_2(x)=x$ for $x\le0$ and
$\widetilde\Gamma_2(t^2)=\widetilde g_2(t)^2$ for $t>0$.
Its second-passage strength is $\widetilde C_{1,2}=C_{1,2}+s$, so
$C_{1,1}+\widetilde C_{1,2}=\widehat m+s=m$.
Reversibility gives its coefficient law and smooth root remainder.
Replace the second passage in the return by $\widetilde\Gamma_2$ and
use the independent translation $\delta$. This recovers the original
return at $(m,\delta)=(\widehat m(\lambda),\widehat\delta(\lambda))$. Joint
smoothness on compact reductions gives the uniform bounds.
\end{proof}

For the estimates below, a family with dependent controls is first treated
in this auxiliary root-profile family and then restricted to
$(m,\delta)=(\widehat m(\lambda),\widehat\delta(\lambda))$.
This convention concerns estimates of return maps; it asserts no field
realization of the auxiliary profiles.

\subsection{Divisions and estimates up to the thresholds}

Where both passages are active, the return admits a division in root
profiles. To cover the thresholds as well, we also estimate the difference
from the aligned return in the coordinate $q$.

\begin{lemma}[Division and estimates up to the thresholds]
\label{lem:exact-value-division}
There are common $r_0>0$, $v_0=1/4$ and a small parameter neighbourhood
in $(m,p)$
such that, for $0<r\le r_0$, $|v|\le v_0$, $q=r^2$ and $\delta=r^2v$,
both passages are active and the return has the expression
\begin{equation}\label{eq:actual-both-active-division}
 Q-q=mr^3+\tfrac34m^2r^4+r^5L_0+\delta rH,
 \qquad H(0,0,0,p)=\tfrac32A(p),
\end{equation}
where $L_0,H$ are restrictions of jointly smooth functions of $(r,m,v,p)$
on the signed box $|r|\le r_0$, $|v|\le v_0$, with the same parameter
neighbourhood. For arbitrary relative sizes and either sign of
$q$ and $\delta$ in the common local domain, including the thresholds,
\begin{equation}\label{eq:quartic-global-value}
 |Q-q-mq_+^{3/2}|
 \le K\{m^2q_+^2+q_+^{5/2}
       +|\delta|\sqrt{q_++|\delta|}+|\delta|q_+^{3/2}\}
\end{equation}
with a common constant.
\end{lemma}

\begin{proof}
Write $g_\nu(t,\lambda)=tT_\nu(t,\lambda)$ with $T_\nu$ jointly smooth
and $T_\nu(0,\lambda)=1$. For signed $r$ and $\delta=r^2v$, put
\[
 U=T_1(r,\lambda)^2,\quad W=U-v,\quad s_2=r\sqrt W,
 \qquad K_0=v+WT_2(s_2,\lambda)^2.
\]
Since $T_1(0,\lambda)=1$, one common reduction gives $T_1^2>3/4$.
For $|v|\le v_0=1/4$ this implies $W=T_1^2-v>1/2$. These expressions are smooth and
$Q(r^2)=r^2K_0$ exactly for $r>0$. Since $K_0(0,m,v,p)=1$,
joint Hadamard division gives $K_0-1=r\mathcal Q$. Aligned cancellation gives
$\mathcal Q(r,m,0,p)=m+\tfrac34m^2r+r^2L_0(r,m,p)$; division in the independent
variable $v$ then gives the exact identity
\begin{equation}\label{eq:quartic-Hadamard}
 \mathcal Q=m+\tfrac34m^2r+r^2L_0+vH(r,m,v,p),
 \qquad H(0,0,0,p)=\tfrac32A(p).
\end{equation}
Indeed, the first root jet at $r=0$ is
$A_1+A_2(1-v)^{3/2}$, whose $v$ derivative at cancellation is
$-3A_2/2=3A(p)/2$. The profile arguments contain $\delta=r^2v$.
Multiplication by $r^3$ proves
\eqref{eq:actual-both-active-division}.

Under~\eqref{eq:rank}, hold $(a,b,m,\xi)$ fixed. For dependent controls,
write $p$ for the original physical parameters and hold $(p,m)$ fixed in
the auxiliary family of Lemma~\ref{lem:auxiliary-controls}, varying only
its auxiliary separation $\delta$. After integration, restrict to the
observed controls as agreed above. Put $y=\Gamma_1(q)-\delta$. The total
derivative is
\begin{equation}\label{eq:quartic-total-detuning}
 Q_\delta=1-\Gamma_{2,q}(y)
 +\Gamma_{2,q}(y)\Gamma_{1,\delta}(q)+\Gamma_{2,\delta}(y).
\end{equation}
The profile parameter derivatives are $O(x_+^{3/2})$. Integrate \eqref{eq:quartic-total-detuning} from zero to $\delta$,
with the indicated quantities held fixed. At an intermediate separation $t$,
$y_+=(\Gamma_1(q,t)-t)_+\le Kq_++|t|$. Joint $C^1$ regularity
and the smooth profiles give, including at active arguments equal to zero,
\[
 |\Gamma_{\nu,q}(x)-1|\le K\sqrt{x_+},\qquad
 |\Gamma_{\nu,t}(x)|\le Kx_+^{3/2}.
\]
Thus $|Q_t|\le K[\sqrt{q_++|t|}+q_+^{3/2}]$, and
\begin{equation}\label{eq:quartic-global-detuning}
 |Q_{m,\delta}-Q_{m,0}|
 \le K|\delta|[\sqrt{q_++|\delta|}+q_+^{3/2}].
\end{equation}
Together with \eqref{eq:passage:aligned-jet} this proves
\eqref{eq:quartic-global-value}.
\end{proof}

\begin{remark}[Checking the hypotheses]\label{cor:field-criterion}
With the prescribed contact geometry, transverse coorientation and
common first-return itinerary, the remaining conditions involve finitely
many jets. The transported strengths are $A_\nu=\tfrac43\mu_\nu(0,0)$;
transport from an initial second-port coordinate multiplies its strength by
$\sqrt{S_1'(\delta)}$, as in Proposition~\ref{prop:field-recognition}.
At the reference, check $A_1=-A_2\ne0$, $R(0)=0$, $R'(0)=1$ and
$c_0=R''(0)/2\ne0$. These conditions give admissibility and the upper
bounds of Theorem~\ref{thm:six-main} and part~1 of
Theorem~\ref{thm:main-saturation-summary}. Attainment additionally uses
\eqref{eq:rank}; the saturation and fourfold-cycle conclusions use the
sign sector specified in Theorems~\ref{thm:main-saturation-summary}
and~\ref{thm:actual-quartic-summary}.
\end{remark}

The curvature of this composition gives the count in the next section, and
the value estimates locate the six-cycle configurations.
